\documentclass[11pt]{article}
\usepackage[margin=2.4cm]{geometry}
\usepackage{fontspec}
\setmainfont{DejaVu Sans}
\usepackage{juliamono}
\usepackage{fvextra}
\usepackage{xcolor}
\usepackage{xurl}
\usepackage{hyperref}
\usepackage{amssymb}
\newcommand{\checkbox}{$\square$}
\newcommand{\answerchoices}{\checkbox\ YES \quad \checkbox\ NO \quad \checkbox\ CANNOT DETERMINE}
\newcommand{\decisionexplanation}{%
  \par\smallskip\textbf{If NO, give the reason for the mismatch (required). If
  CANNOT DETERMINE, state what context or evidence is missing:}\par
  \rule{\linewidth}{0.4pt}\par\vspace{7pt}\rule{\linewidth}{0.4pt}\par}
\setlength{\parindent}{0pt}\setlength{\parskip}{5pt}
\setlength{\emergencystretch}{2em}
\DefineVerbatimEnvironment{Vsrc}{Verbatim}{fontsize=\small,frame=single,
  breaklines=true,breakanywhere=true}
\DefineVerbatimEnvironment{Vlean}{Verbatim}{fontsize=\scriptsize,frame=single,
  numbers=left,numbersep=5pt,breaklines=true,breakanywhere=true}
\hypersetup{colorlinks=true,urlcolor=blue}

\begin{document}
\begin{center}{\LARGE Proof Verification Document}\\[2pt]{\large reddit\_mathematics\_i3ob36}\\[2pt]{\small Reddit proof-check posts (newly confirmed)}\end{center}
\begin{center}\fcolorbox{blue!60!black}{blue!4}{\parbox{0.92\linewidth}{\textbf{Make three independent decisions.} Follow the evidence chain: original source $\longrightarrow$ challenged claim $\longrightarrow$ proposed counterexample $\longrightarrow$ Lean theorem. A record is valid only when all three links are correct. Complete the human sign-off before reading the automated review appendix; those earlier results are background, not an answer key.}}\end{center}
\textbf{Use only this document, except for the public source link.} Everything needed for Questions 2 and 3 is included here. For Question 1, open the public source link when one is shown. If there is no public link, the complete stored problem and solution are included here. If the information is still not enough, choose CANNOT DETERMINE.
For each question, choose \textbf{YES}, \textbf{NO}, or \textbf{CANNOT DETERMINE}. Choose CANNOT DETERMINE if the document does not provide enough information for you to decide; do not mark YES if you are uncertain. If you find a specific mismatch, choose NO and explain it in the space below that question. Do not rewrite the claim, add assumptions, or substitute a different counterexample.
\textbf{How copied text is shown.} The problem, quoted passage, and proposed counterexample come from the stored record. The Lean theorem statement and Lean files are copied exactly from the certificate. Some long lines continue on the next printed line; this is only page formatting and does not change the copied text. Line numbers are added only to help you find passages. If the stored record contains a character that cannot be printed, the document shows \texttt{[U+XXXX CONTROL]} in its place.
\begin{description}
\item[Challenged claim.] The exact statement from the source that this record says is false.
\item[Proposed counterexample.] A concrete case that should satisfy the claim's assumptions but make its conclusion false.
\item[Lean certificate.] A Lean file proving a formal theorem. Successful compilation proves the theorem as written, not that it is the right translation of the source.
\end{description}
\section*{Question 1 --- Does this proof verification document accurately represent the source?}
\textbf{What to decide.} Compare the challenged passage below with the original source---either the public link or the complete source text in this document. Answer YES only if they make the same claim and include the same assumptions or conditions. If this document changes the meaning, leaves out a condition, or points to a different passage, answer NO.
\subsection*{Original source}
\textbf{Source record:}
\begin{Vsrc}
qedbench_certificate_availability_triage/data/records_reddit_exhaustive.jsonl
\end{Vsrc}
\textbf{Public source:}\par
\url{https://www.reddit.com/r/mathematics/comments/i3ob36/interesting_proof_by_induction_example/}\par
\subsection*{Problem or question from the source (verbatim)}
\begin{Vsrc}
Interesting proof by Induction example
\end{Vsrc}
\subsection*{Challenged passage recorded from the source (stored text)}
\begin{Vsrc}
If n+1 wasn’t interesting then it would be the smallest non interesting natural number
\end{Vsrc}
\subsection*{Surrounding source lines (stored text with added line labels)}
The labels refer to lines in the solution text stored with the source record; the source-line content is not transliterated.
\begin{Vsrc}
0001 | All natural numbers are interesting.
0002 | 
0003 | We start with the base where n=1. 1 is interesting for many reasons but for example because it is the smallest natural number.
0004 | 
0005 | Then we have the inductive assumption. Assume that n is interesting.
0006 | 
0007 | Now we need to show that n+1 is interesting. If n+1 wasn’t interesting then it would be the smallest non interesting natural number, however this would make it an interesting number, hence n+1 is interesting.
0008 | 
0009 | QED
\end{Vsrc}
\textbf{A YES means:} the proof verification document faithfully represents both the source context and the challenged claim, without dropping a condition or changing its meaning.
\section*{Question 2 --- Does the proposed counterexample disprove that claim?}
\textbf{Claim to test (repeated for convenience).}
\begin{Vsrc}
If n+1 wasn’t interesting then it would be the smallest non interesting natural number
\end{Vsrc}
\subsection*{How the record says the argument uses this claim}
This is the record's description, not an established fact; check it against the source.
\begin{Vsrc}
Assuming only that n is interesting, if n+1 is non-interesting, then n+1 is the smallest non-interesting natural number.
\end{Vsrc}
\subsection*{Proposed counterexample (exact values or construction)}
\begin{Vsrc}
Let natural numbers be {0,1,2,...}, define “interesting” by I={3}, and take n=3. Then n is interesting and n+1=4 is non-interesting. But 1 is also non-interesting and 1<4, so 4 is not the smallest non-interesting natural number.
\end{Vsrc}
\subsection*{Why the record says the counterexample works}
Check this explanation independently; it is evidence to assess, not an answer key.
\begin{Vsrc}
The hypotheses constrain only n and n+1; they impose no condition on smaller naturals. Thus the claimed minimality does not follow.
\end{Vsrc}
\subsection*{How to decide}
\begin{enumerate}
\item List every assumption that applies to the challenged claim.
\item Check that the proposed example satisfies every one of those assumptions.
\item Recompute the claimed outcome and the actual outcome using exact reasoning.
\item Confirm that the conflict is with the challenged claim itself, not an unrelated sentence.
\end{enumerate}
\textbf{A YES means:} the example is fully specified, satisfies all relevant assumptions, and makes the challenged claim false. If a necessary value, condition, or calculation is missing, choose CANNOT DETERMINE.
\section*{Question 3 --- Does the final Lean theorem statement verify this same counterexample?}
\textbf{What to decide.} A Lean file can compile while proving the wrong theorem. Compare these three links independently:
\begin{enumerate}
\item The Lean definitions represent the same mathematical objects, assumptions, and operations as the challenged claim.
\item The Lean construction uses the same concrete values or example shown in Question 2.
\item The final theorem statement asserts that the challenged claim fails on that example, not merely a weaker or nearby statement.
\end{enumerate}
\subsection*{Selected final Lean theorem statement (verbatim UTF-8 excerpt)}
This is an exact contiguous slice of the certificate file, from the theorem keyword through the original \texttt{:=} token where its proof begins. It is shown first so the formal target is not buried in the complete file.
\begin{Vlean}
theorem refutes_source_claim : ¬ SourceClaim :=
\end{Vlean}
\subsection*{Project-local imported Lean files (complete, verbatim UTF-8)}
This section follows the certificate's import commands recursively and prints every repository-owned Lean source file on which the final theorem statement depends. Each file is copied in full, including its definitions and controls. External Mathlib, Std, and Lean Init source is not duplicated.
No project-local imported Lean files were found. External Mathlib, Std, and Lean Init source is not duplicated in this proof verification document.
\subsection*{Complete counterexample Lean file (verbatim UTF-8)}
The complete Lean file is included below so you can check every definition and supporting result used by the final theorem statement. Use it together with the project-local imported Lean files above.
\textbf{Lean certificate file:}
\begin{Vsrc}
partB_feasibility/reddit_certificates/reddit_mathematics_i3ob36.lean
\end{Vsrc}
\begin{Vlean}
/-!
Record: `reddit_mathematics_i3ob36`

Public source: https://www.reddit.com/r/mathematics/comments/i3ob36/interesting_proof_by_induction_example/

Exact target scope: the local induction-step assertion that, for every
interestingness predicate `I` on the natural numbers and every natural `n`,
the two premises `I n` and `¬ I (n + 1)` alone imply that `n + 1` is the
smallest non-interesting natural.  The base case is a separate obligation and
is not modeled as a premise of this local implication.

Modeling judgments: natural numbers are Lean's actual `Nat` values
`0, 1, 2, ...`; interestingness is an arbitrary predicate `Nat → Prop`;
addition and strict order are Lean's ordinary natural-number operations; and
`SmallestNonInteresting I m` means both `¬ I m` and that every natural `k < m`
is interesting.  The source implication retains both universal quantifiers
and both load-bearing hypotheses.  No finite proxy is introduced.

Concrete witness: take `I = {3}` and `n = 3`.  Then `I 3`, while `I 4` is
false, but `I 1` is also false and `1 < 4`; hence `4` is not the smallest
non-interesting natural.  The independently verified negative control uses
the same `n = 3` and `I = {0,1,2,3}`, encoded by `I k ↔ k < 4`: its premises
hold and `4` genuinely is the smallest non-interesting natural.  Thus the
control detects a vacuous or uniformly-false encoding and checks the omitted
smaller-number condition on which minimality depends.

Route review: this is a self-contained Lean certificate because the claim and
both models use only genuine natural numbers, ordinary addition and order,
predicates, and first-order quantification.  The implementation deliberately
uses zero imports and kernel `decide` only for exposed closed numeral facts.
-/

namespace RedditMathematicsI3ob36

def SmallestNonInteresting (I : Nat → Prop) (m : Nat) : Prop :=
  ¬ I m ∧ ∀ k : Nat, k < m → I k

def SourceClaim : Prop :=
  ∀ (I : Nat → Prop) (n : Nat),
    I n → ¬ I (n + 1) → SmallestNonInteresting I (n + 1)

def witnessI (k : Nat) : Prop :=
  k = 3

def controlI (k : Nat) : Prop :=
  k < 4

theorem witness_three_is_interesting : witnessI 3 := by
  rfl

theorem witness_four_is_not_interesting : ¬ witnessI (3 + 1) := by
  change ¬ ((4 : Nat) = 3)
  decide

theorem witness_one_is_not_interesting : ¬ witnessI 1 := by
  change ¬ ((1 : Nat) = 3)
  decide

theorem witness_one_is_smaller : 1 < 3 + 1 := by
  decide

theorem witness_four_is_not_smallest :
    ¬ SmallestNonInteresting witnessI (3 + 1) := by
  intro hSmallest
  exact witness_one_is_not_interesting
    (hSmallest.2 1 witness_one_is_smaller)

/-- Load-bearing refutation of the literal, fully quantified source claim. -/
theorem refutes_source_claim : ¬ SourceClaim := by
  intro hClaim
  exact witness_four_is_not_smallest
    (hClaim witnessI 3 witness_three_is_interesting
      witness_four_is_not_interesting)

theorem control_three_is_interesting : controlI 3 := by
  change (3 : Nat) < 4
  decide

theorem control_four_is_not_interesting : ¬ controlI (3 + 1) := by
  change ¬ ((4 : Nat) < 4)
  decide

/-- Meaningful negative control: under `I k ↔ k < 4`, the same local premises
hold and the asserted minimality conclusion is true. -/
theorem negative_control_four_is_genuinely_smallest :
    controlI 3 ∧
    ¬ controlI (3 + 1) ∧
    SmallestNonInteresting controlI (3 + 1) := by
  constructor
  · exact control_three_is_interesting
  constructor
  · exact control_four_is_not_interesting
  constructor
  · exact control_four_is_not_interesting
  · intro k hk
    exact hk

end RedditMathematicsI3ob36

\end{Vlean}
\subsection*{Proof-safety information}
The automated scan found no unfinished proofs, new axiom declarations, unsafe declarations, external code hooks, or uses of Lean's compiled decision procedure. This scan concerns proof construction only; it does not establish that the theorem matches the source claim.
\textit{Technical scan counts:} native\_decide=0, sorry=0, axiom\_decl=0, unsafe=0, extern=0. \textit{Technical tier:} KERNEL-ONLY.
\textbf{A YES means:} Lean represents the same objects and assumptions, uses the same proposed counterexample, and the final theorem statement asserts the same contradiction. Successful compilation by itself is not enough.
\clearpage
\section*{Human sign-off}
Complete this block from your own review of Questions 1--3. Do not use the automated appendix below as an answer key.\par
\textbf{Reviewer name:} \rule{9cm}{0.4pt}\par
\textbf{Date:} \rule{4cm}{0.4pt}\par
\medskip\textbf{1. Does the proof verification document accurately represent the source and challenged claim?}\par
\answerchoices\par
\decisionexplanation
\medskip\textbf{2. Does the proposed counterexample satisfy the assumptions and make that claim false?}\par
\answerchoices\par
\decisionexplanation
\medskip\textbf{3. Does the final Lean theorem statement represent and verify that same counterexample?}\par
\answerchoices\par
\decisionexplanation
\medskip\textbf{Overall decision:}\par
\checkbox\ VALIDATED (all three answers are YES)\par
\checkbox\ REJECTED (at least one answer is NO)\par
\checkbox\ NEEDS ANOTHER REVIEWER (at least one CANNOT DETERMINE and no NO)\par
\medskip\textbf{Additional comments.} Use this space for any extra source, counterexample, calculation, or Lean details that support your decisions.\par\rule{\linewidth}{0.4pt}\par\vspace{7pt}\rule{\linewidth}{0.4pt}
\clearpage
\section*{Appendix --- Prior automated review (background only)}
\textbf{These results are not an answer key.} They come from the earlier Opus/Fable verification campaign and are printed only so an auditor can see what was previously checked. The human decisions above must be made independently. Raw status codes are retained in small type for traceability.
\textbf{Overall automated disposition:} Cleared by automated screening; an independent search also found a counterexample.
\par\textit{Audit code:} \nolinkurl{PASS_CORROBORATED}
\subsection*{Independent search with no access to the stored example}
\textbf{Result:} An independent search found a counterexample without seeing the stored example.
\par\textit{Audit code:} \nolinkurl{WITNESS_FOUND}
\par\textbf{Checker's explanation:}
\begin{Vsrc}
Universe {0,1,2}; predicate I(0)=False, I(1)=True, I(2)=False; induction index n=1. Under the WEAK (ordinary) induction hypothesis both premises of the quoted line hold -- I(n)=I(1)=True and ~I(n+1)=~I(2)=True -- yet the conclusion fails: the claim asserts min{m : ~I(m)} = n+1 = 2, while in fact min{m : ~I(m)} = 0. So n+1 is non-interesting without being the SMALLEST non-interesting natural number, and the inference 'if n+1 wasn't interesting then it would be the smallest non-interesting natural number' is invalid as stated. (Shortened here; the complete note remains in the machine-check JSON.)
\end{Vsrc}
\par\textbf{Cases checked:}
\begin{Vsrc}
Exhaustive over: start in {0,1} (both 'natural number' conventions); universe {start..N} for every N with 0 <= N-start <= 15; every truth assignment I : {start..N} -> {True,False} (all 2^(N-start+1) predicates, enumerated by itertools.product in fixed lexicographic order); every induction index n with start <= n <= N-1 so that n+1 lies in the universe. (Shortened here; the complete note remains in the machine-check JSON.)
\end{Vsrc}
\subsection*{Independent re-check of the proposed counterexample}
\textbf{Result:} A separate checker recomputed the proposed counterexample and confirmed it.
\par\textit{Audit code:} \nolinkurl{CONFIRMED}
\par\textbf{Checker's note:}
\begin{Vsrc}
Independently recomputed with exact integer arithmetic, stdlib only, no floats or randomness in any decision. Under the witness (I={3}, n=3): n=3 is interesting and n+1=4 is non-interesting, so the quoted step's hypotheses hold; but 0, 1 and 2 are all non-interesting and less than 4, so 4 is not the smallest non-interesting natural. The consequent fails, hence the claim fails on the witness. The witness text cites 1 as the smaller non-interesting number; (Shortened here; the complete note remains in the machine-check JSON.)
\end{Vsrc}
\subsection*{Independent review of the source-to-Lean correspondence}
\textbf{Result:} A separate reviewer judged that Lean proves the same failure using the same example.
\par\textit{Audit code:} \nolinkurl{MATCHES}
\par\textbf{Theorem summary from that reviewer:}
\begin{Vsrc}
The file defines SmallestNonInteresting I m := (not I m and forall k < m, I k) and SourceClaim := forall (I : Nat -> Prop) (n : Nat), I n -> not I (n+1) -> SmallestNonInteresting I (n+1). The load-bearing theorem refutes_source_claim : not SourceClaim discharges the claim by instantiating it at witnessI k := (k = 3) and n = 3: witness_three_is_interesting (witnessI 3), witness_four_is_not_interesting (not witnessI (3+1)), and witness_four_is_not_smallest (not SmallestNonInteresting witnessI. (Shortened here; the complete note remains in the machine-check JSON.)
\end{Vsrc}
\par\textbf{Difference noted by that reviewer:}
\begin{Vsrc}
None material. (1) Definitions: 'smallest non-interesting natural' is encoded as not I m and forall k < m, I k, which is logically equivalent to 'not I m and forall k, not I k -> m <= k' over Nat -- neither weakened nor strengthened; interestingness is an arbitrary Nat -> Prop, naturals are Lean's Nat (including 0), and both universal quantifiers plus both load-bearing hypotheses (I n, not I (n+1)) are retained; no finite proxy. (Shortened here; the complete note remains in the machine-check JSON.)
\end{Vsrc}
\subsection*{Follow-up review}
No follow-up review was needed or recorded.
\textit{This appendix is only background. You do not need any separate machine-output files to answer the three questions.}
\end{document}